% Step 11: Circle reduction of W+ anisotropy and the fiber survival dichotomy
% S0910A03 Agent2 (Constructor) — Explore
%
% This note develops the static circle reduction of W+ for the
% negative-branch calibrated sequence, corrects the initial formula
% (following TTCS verification), and establishes the key dichotomy
% connecting W+ anisotropy to fixed-scale fiber survival.
%
% Supersedes: the uncorrected formulas in circle_reduction_weyl.py and
% verify_static_incompatibility.py.  Corrected script: corrected_circle_reduction.py.
%
% STEP 12 REPAIR: the static formula in Sections 1--2 remains valid, as does
% the spherical cohomogeneity-one calibrated exclusion.  The claimed general
% static rigidity, the fiber-survival/high-twist dichotomy, and the inference
% from toric orthogonal transitivity to a static chosen circle are withdrawn.
% Small ell^2|F|^2 does not control ell*nabla F or (nabla ell)*F, and quotient
% geometry was not derived.  See workspace/notes/step12_reduced_equations_accessibility.md.

\section{Static circle reduction of $W^+$: corrected formula}

Let $(M^4,g)$ be an oriented Riemannian four-manifold with a static free
circle action, so $g=q+\ell^2d\psi^2$ where $q=g_B$ is the base metric on
$B^3$ and $\ell:B\to\R_{>0}$ is the orbit-length function.  If $\Ric_g=-g$,
the warped-product Ricci equations give
\begin{equation}\label{eq:static-einstein}
 \Ric_B=h-g_B,\qquad \Delta_B\ell=\ell,\qquad
 h\coloneqq\ell^{-1}\operatorname{Hess}_B\ell,\quad \tr h=1.
\end{equation}
In the adapted orthonormal frame $e_1,e_2,e_3$ for $B$ and $e_4=\ell^{-1}\partial_\psi$,
the self-dual 2-forms are
\[
 \omega_1=\frac{e^{12}+e^{34}}{\sqrt2},\qquad
 \omega_2=\frac{e^{13}-e^{24}}{\sqrt2},\qquad
 \omega_3=\frac{e^{14}+e^{23}}{\sqrt2}.
\]
The static warped-product curvature gives $R_{abcd}=R^B_{abcd}$,
$R_{a4b4}=-h_{ab}$, $R_{abc4}=0$.  In three dimensions,
$R^B_{abcd}$ is determined by $\Ric_B=h-I$ and $s_B=-2$.

\medskip
\noindent\textbf{Proposition (Corrected static reduction).}
\emph{The curvature operator on $\Lambda^2_+$ in the Besse convention is}
\begin{equation}\label{eq:rhat-static}
 \widehat R_+=
 \begin{pmatrix}
  -h_{33}&h_{23}&-h_{13}\\
  h_{23}&-h_{22}&h_{12}\\
  -h_{13}&h_{12}&-h_{11}
 \end{pmatrix}.
\end{equation}
\emph{Consequently,}
\begin{equation}\label{eq:spec-rhat}
 \boxed{\operatorname{Spec}(\widehat R_+)=-\operatorname{Spec}(h),}
\end{equation}
\emph{and}
\begin{equation}\label{eq:wplus-static}
 W^+=\widehat R_++\frac13I,
 \qquad
 \boxed{\operatorname{Spec}(W^+)
  =\left\{\frac13-\eta_i\right\}_{i=1,2,3},}
\end{equation}
\emph{where $\eta_1,\eta_2,\eta_3$ are the eigenvalues of $h$.}

\begin{proof}
For the diagonal entries, use $\tr h=1$ and the three-dimensional identity
$K_{ab}=(h_a+h_b-h_c-1)/\ldots$ to simplify:
\[
 (\widehat R_+)_{11}=\frac12(K_{12}-h_{33})
 =\frac{h_{11}+h_{22}-3h_{33}-1}{4}=-h_{33},
\]
where the last step uses $h_{11}+h_{22}=1-h_{33}$.  Similarly for the other
diagonal entries.  Off-diagonal entries use $R^B_{1213}=\Ric^B_{23}=h_{23}$
and $R_{a4b4}=-h_{ab}$.  The eigenvalue statement follows because the matrix
\eqref{eq:rhat-static} is conjugate to $-h$ by the index permutation
$1\leftrightarrow3$ (a diagonal orthogonal transformation).
\end{proof}

\noindent\textbf{Calibrated Hessian spectrum.}
The calibrated self-dual Weyl spectrum
$(-\frac16,\frac1{12},\frac1{12})$ requires
\begin{equation}\label{eq:h-calibrated}
 \boxed{\operatorname{Spec}(h)=\left(\frac12,\frac14,\frac14\right),
 \qquad \tr h=1.}
\end{equation}
The simple $W^+$ eigenvalue $-1/6$ corresponds to the large Hessian
eigenvalue $1/2$; the double eigenvalue $1/12$ corresponds to the
smaller pair $1/4$.  This means the warping function $\ell$ has
\emph{stronger} curvature in the direction singled out by the simple
$W^+$ eigenline.

\medskip
\noindent\textbf{Verification.}
For constant curvature $K=-1/3$: $h=\frac13I$, $\widehat R_+=-\frac13I$,
$W^+=0$.  Script \texttt{corrected\_circle\_reduction.py} confirms both
the isotropic and calibrated cases.

\section{Static cohomogeneity-one incompatibility}

For $g=dt^2+a(t)^2g_{S^2}+\ell(t)^2d\psi^2$ with $\Ric=-g$, let
$u=\ell''/\ell$ and $v=a'\ell'/(a\ell)$.  Then $h=\diag(u,v,v)$ and
\[
 \Spec(W^+)=\left(\frac13-v,\frac13-v,\frac13-u\right).
\]
The calibrated target gives $u=\frac12$, $v=\frac14$.

\medskip
\noindent\textbf{Proposition (Static ODE incompatibility).}
\emph{There is no smooth cohomogeneity-one static circle Einstein metric
with constant calibrated $W^+$ spectrum.}

\begin{proof}
From the radial Einstein equation:
\[
 a''=\frac{a(1-u)}{2}=\frac{a}{4}.
\]
From the angular Einstein equation together with the calibrated $v$:
\[
 (a')^2=1+\frac{a^2}{2}.
\]
Differentiating: $a''=a/2$, contradicting $a''=a/4$.  The relative gap is
$|a/2-a/4|/(a/2)=1/2$.
\end{proof}

\noindent\textbf{Scope of the ODE calculation.}
If constant $u,v$ are allowed to float \emph{within this spherical
cohomogeneity-one ansatz}, the same elimination gives $u=v$.  Combined with
$u+2v=1$, this forces $u=v=\frac13$ and hence $W^+=0$.  This is not a
classification of general static Einstein circle metrics.  For example,
$\mathbb H^2\times\mathbb H^2$, using polar circles in the first factor, has
$\operatorname{Spec}(h)=(1,0,0)$ and constant
$\operatorname{Spec}(W^+)=(-2/3,1/3,1/3)$ up to orientation.

\section{Twisted reduction and sign correction}

For a general Kaluza--Klein metric $g=q+\ell^2\theta^2$, $\theta=d\psi+A$,
$F=dA$, the vertical Einstein equation becomes
\[
 -\frac{\Delta_B\ell}{\ell}+\frac{\ell^2}{4}F_{ab}F^{ab}=-1,
\]
giving
\begin{equation}\label{eq:tr-h-twisted}
 \boxed{\tr h=1+\frac{\ell^2}{4}F_{ab}F^{ab}\ge1,}
\end{equation}
with equality only when $F=0$.  This corrects the earlier sign error;
twist \emph{increases} $\tr h$ above the static value $1$.

The horizontal Einstein equation is modified to
\[
 \Ric_B=h+\frac{\ell^2}{2}(F^2)_{\rm sym}-g_B,
\]
and the full curvature operator on $\Lambda^2_+$ acquires additional terms:
\[
 \widehat R_+=\widehat R_+^{\rm static}+\widehat R_+^{\rm twist}.
\]
For the Taub--NUT--AdS family,
$\widehat R_+^{\rm static}+\widehat R_+^{\rm twist}=-\frac13I$
(i.e., $W^+=0$ in the half-flat orientation).

\section{Near-NUT behavior: corrected analysis}

A critical correction from the TTCS verification: near a smooth NUT
where $\ell\to0$, the quotient-base geometry degenerates.  Setting
$d=r-\nu$ (distance from the NUT in the Taub--NUT--AdS radial coordinate),
\[
 \ell^2\sim2\nu d,\qquad
 h_{\rm ang}\sim\frac{1}{8\nu d}\to\infty,\qquad
 \ell^2|F|^2\sim\frac{1}{2\nu d}\to\infty.
\]
Both the static Hessian contribution and the twist correction diverge
individually while canceling in the total (smooth) curvature.

\medskip
\noindent\textbf{Consequence.}
The static formula $W^+\approx\frac13I-h$ does \emph{not} dominate near
a smooth NUT.  The net $W^+$ depends on the precise balance between
divergent Hessian and twist terms.  In particular, it is \emph{not}
correct to conclude that ``$W^+\to0$ near any smooth NUT.''  A
NUT with appropriate twist could in principle produce any $W^+$
eigenvalue structure.

This removes one initially proposed mechanism (NUT forces isotropic $W^+$,
creating calibrated deficit) and replaces it with the more precise dichotomy
below.

\section{Withdrawn fiber-survival/high-twist dichotomy}

The following division was proposed by the TTCS generator but is \emph{not
established}; it is retained only as a record of the superseded proposal.

\medskip
\noindent\textbf{Superseded proposed division.}
The following alternatives were conjectured for the calibrated good set, but
the asserted implication in the second item was not proved:

\begin{enumerate}
\item[\textbf{(S)}] \textbf{Fiber survival.}
\emph{The circle fiber class survives at a fixed buffer scale
$\rho\ge\rho_0(K)>0$: a buffered normal cover at that scale unwraps
the fiber, the cover is noncollapsed, and the mean-transport carries
the calibrated deficit.  This produces the forbidden calibrated Einstein
germ (Steps~5--6), giving a contradiction.}

\item[\textbf{(D)}] \textbf{Fiber death with high twist.}
\emph{The fiber contracts within every fixed buffer scale.  The static
incompatibility forces}
\[
 \ell^2|F|^2=\tau\ge c_Q>0
\]
\emph{on a definite fraction of the region; otherwise the leading-order
$W^+$ would be forced toward $W^+=0$ (static isotropic), contradicting
the calibrated deficit bound.}
\end{enumerate}

The fiber-length function $\ell(p)$ can be arbitrarily small at a good
point (a conformal rescaling $\ell\mapsto\varepsilon\ell$ does not change
$h=\Hess\ell/\ell$).  The curvature mean bounds $h$, not $\ell$.

In the non-degenerate case (\textbf{S}), the Harnack estimate for the
positive eigenfunction equation $\Delta_B\ell=(1+\tau/4)\ell$ gives
\[
 \sup_{B_{R/2}}\ell\le C_H\inf_{B_{R/2}}\ell,\qquad
 |\nabla\log\ell|\le C_H/R,
\]
provided the base geometry, twist energy $\tau$, and base harmonic
radius are controlled.  Under these conditions, a positive fiber length
at a good point remains positive throughout a fixed-radius neighborhood,
supplying the buffer domain for the covering chart.

Case \textbf{(D)} requires a diffuse mechanism:
$\ell^2|F|^2\ge c>0$ must persist on a positive-volume subset to maintain
the calibrated anisotropy.  This ``scaled twist energy'' acts as a source
of non-basic curvature.

\section{Relation to Pontryagin and accessibility}

The Strategist's correction clarifies that the positive lower bound on
$|\Xi|^2=|\iota_VF_{A_+}|^2$ (from the calibrated Pontryagin density,
Step~10) is \emph{not discriminating}: for any Einstein metric with
$\Ric=-g$, the Levi-Civita connection on $\Lambda^2_+$ has self-dual
curvature, giving $|\iota_VF|^2=|F|^2/2$ automatically.  The isotropic
scalar contribution $s^2/24$ provides $88.9\%$ of the calibrated
characteristic density; the anisotropic $2|W^+|^2$ term adds only
$11.1\%$.

The \emph{distinguishing} information is the eigenvalue structure:
the calibrated curvature operator on $\Lambda^2_+$ has spectrum
$(-1/2,-1/4,-1/4)$, while the isotropic (hyperbolic) spectrum is
$(-1/3,-1/3,-1/3)$.  The anisotropy
\[
 |W^+|^2=1/24
\]
is the correct measure of deviation from the benign Einstein background,
not the total characteristic density.

\section{Peer results and integration}

S0910A01 (Step~12): independently verified the Taub--NUT--AdS multiscale
computation and agrees that the target is multiscale rank recovery, not
NUT charge counting.  Their rank-one Cartan structure equations support
the diffuse mixed-curvature picture.

S0910A02 (Step~12): developed toric determinant calculations for $T^2$-invariant
metrics.  Key finding: $\tr_h\Ric^V=-\Delta_Bw/w+|F|_h^2/2$ with
$w=\sqrt{\det h}$, leading to $\Delta_Bw=(|F|^2/2-2\lambda)w$.
For $\lambda=-1$: $\Delta_Bw=(|F|^2/2+2)w$, giving strong subharmonicity
of the orbit-volume function.  Their orthogonal transitivity lemma
(conserved twist constants vanish at degeneration edges) is relevant to
our dichotomy: if twist vanishes near bolts/edges, the remaining collapse
is static and incompatible with the calibrated spectrum.

\section{Current research position (repaired in Step 12)}

\textbf{What is established:}
\begin{itemize}
\item The corrected static circle reduction formula
$\Spec(W^+)=\{1/3-\eta_i\}$, verified against constant curvature,
$\mathbb H^2\times\mathbb H^2$, and Taub--NUT--AdS.
\item The calibrated spectrum is incompatible with the displayed spherical
cohomogeneity-one static ansatz.  General calibrated exclusion comes instead
from the Step 6 projector/Derdzi\'nski Einstein-germ argument.
\item Buffered fiber survival plus a supplied controlled cover chart leads to
the forbidden germ.  No converse from failure of accessibility to high
$\ell^2|F|^2$ has been proved.
\item The Pontryagin density is not discriminating; the $W^+$ anisotropy is.
\end{itemize}

\textbf{What remains open:}
\begin{itemize}
\item Can derivative twisting $\ell\nabla F+(\nabla\ell)F$ or quotient
geometry degeneration prevent every fixed-scale buffered chart while the
calibrated profile and curvature means remain controlled?
\item Can the weighted quotient structure inherited from four-dimensional
Ricci bounds be upgraded, together with the complete reduced equations, to
fixed-scale cover reconstruction?  The earlier Harnack argument assumed the
ordinary base control it was meant to obtain.
\item The positive branch, higher dimensions, and Lean formalization.
\end{itemize}

\section{Files and references}
\begin{itemize}
\item Corrected script: \texttt{workspace/scripts/corrected\_circle\_reduction.py}
\item TTCS generator: \texttt{trace/ttcs/results/step11\_ttcs01\_generator1\_integration.md}
\item TTCS verifier: \texttt{trace/ttcs/results/step11\_ttcs01\_verifier1\_integration.md}
\item Previous notes: Steps~6, 9, 10 (local obstruction, Weitzenb\"ock, Pontryagin)
\end{itemize}
